% !TeX root = ../../Book1.tex
% 纲要标题：### A.4 拓扑预备知识的阅读建议
% 纲要位置：计划纲要.md，第 805 行。
% BEGIN APPENDIX OUTLINE SYNC
% <!-- 短节：不设小节 -->
% END APPENDIX OUTLINE SYNC

\section{拓扑预备知识的阅读建议}
\label{b1:sec:A04}

群、环和有限域的基础章节不需要完整的拓扑课程。商对象的底集合可由\secref{b1:sec:A02}核对；遇到一般极大理想、超越基或域嵌入延拓时，查阅\secref{b1:sec:A01}，重点检查偏序、初始对象、链上界及极大性如何推出目标。这些章节所需的代数证明均在正文，附录只把共同的集合论语言放在一处。

\ProseParagraphBreak
阅读\chapref{b1:ch:22}时，可以先沿正文的有限 Galois 层进入逆极限，再按需要回看\secref{b1:sec:A03}。必须掌握的是开集原像、有限坐标邻域、闭集的有限交性质和紧空间到 Hausdorff 空间的连续双射判据。度量化、序列紧性和一般积紧性理论都不是该章的额外前置。本附录也没有证明“每个紧 Hausdorff 全不连通空间都能写成有限空间逆极限”的反向刻画，不能仅凭这里的一向结论使用它。

\ProseParagraphBreak
后续卷中的素谱通常只具拟紧性，未必 Hausdorff；完备化中的逆极限还涉及环、模和正合性。它们需要分别增加代数或拓扑假设，不能把有限离散集合的结论原样照搬。例如各层非空但无限时逆极限可以为空，连续双射若缺少紧源或 Hausdorff 目标也未必为同胚。查阅其他教材时，应先核对“紧”的约定、过渡映射的方向以及是否假定满射，再比较定理的名字。
% ---
